\startchapter[title={Revisão}]
	A indústria de biotecnologia é responsável por produções alternativas de compostos usuais na nossa sociedade, utilizando organismos vivos para produzir alimentos, biocombustíveis \cite[authoryear][ref::nielsen_0], fármacos \cite[authoryear][ref::zhang_0] e compostos químicos \cite[authoryear][ref::meadows_0]. No balanço de sustentabilidade, estas rotas alternativas, utilizando organismos modificados ou não, permitem uma aproximação orgânica e amigável com o meio ambiente para os processos produtivos envolvidos, promovendo uma transição para uma economia verde, ou seja, reduzindo o saldo de carbono emitido \cite[authoryear][ref::naseri_0]. Estas alternativas não são à princípio uma digressão em produtividade nas cadeias produtivas mencionadas em prol da sustentabilidade, mas sim, puramente uma evolução em eficiência na alocação de espaço físico, consumo de insumos e geração de subprodutos \cite[authoryear][ref::metcalfe_0]. A sustentabilidade é mais um ponto positivo que pesa a favor de implementações em grande escala de processos que utilizem técnicas em biotecnologia em suas rotas produtivas.

	O que sustenta a reprodutibilidade dos processos bioquímicos industriais é a possibilidade de monitorar a qualidade destes processos com sensores físicos. Estes são de suma importancia para que os operadores ou os controladores tenham uma medida aproximada de como o processo está ocorrendo e se alguma contra-medida precisa ser tomada para contornar problemas diversos. Existem alternativas baseadas em {\emph software} para mensurar valores de variáveis em processos bioquímicos, estes precisam ser desenvolvidos programáticamente utilizando técnicas de modelagem matemática.

	Os modelos matemáticos estudados para o presente trabalho para formular um sensor baseado em {\emph software} especializado em medir concentração de hidrogênio em um fermentador de hidrogênio são apresentados a seguir.

	% \startsection[title={Biorreatores}]
  Biorreatores fornecem ambientes favoráveis para o crescimento de culturas de microrganismos e podem ser otimizados para a produção de determinado produto bioquímico, desde que haja acesso a mecanismos de controle das variáveis físicas envolvidas, o que geralmente é oferecido por reatores projetados especificamente para essa finalidade, contrariamente aos bioreatores de ocorrencia natural, onde a manipulaçao consistente de variaveis nao acontece com facilidade ou é completamente mitigada - a exemplo, uma poça de agua a ceu aberto com atividade microbiana em que se deseja controlar a temperatura. A necessidade de um projeto de biorreator decorre da responsabilidade de garantir a reprodutibilidade dos processos bioquímicos desenvolvidos ao longo de uma cadeia produtiva. A possibilidade de controlar algumas variáveis \emdash ou até mesmo todas aquelas avaliadas como impactantes para a conversão final \emdash é uma característica indispensável para aplicações industriais de rotas produtivas desenvolvidas em laboratório. Sendo que uma vez determinadas as condições de operação na etapa experimental, basta ajustar-se as variáveis extensivas do processo para uma escala industrial.

  Segundo \cite[authoryears][ref::erickson_0], existem diversas formas de classificar biorreatores, seja pelo tipo de alimentação ou pelos produtos finais obtidos. Existem três tipos principais de alimentação: batelada, alimentação alternada e alimentação contínua. A alimentação em batelada ocorre quando o substrato é introduzido no sistema apenas no início do processo, sem manipulações externas posteriores na quantidade de fonte de carbono. A alimentação alternada consiste em uma estratégia de distribuição da fonte de carbono seguindo um cronograma, diferenciando-se da primeira pela possibilidade de introdução de substrato durante o experimento. Por fim, a alimentação contínua caracteriza-se pela introdução constante de alimento ao longo de todo o processo, representando uma tentativa de reproduzir um sistema estacionário a partir da manutenção da concentração de substrato no biorreator. A \in{Figura}[fig:simple_bioreactor] mostra um examplo de volume de controle que pode ser pensado como um bioreator.

  \startplacefigure[title={Ilustração de biorreator genérico}, reference={fig:simple_bioreactor}]
    \externalfigure[image/simple_bioreactor.png][height=200pt]
  \stopplacefigure

  O bioreator representado na \in{Figura}[fig:simple_bioreactor] tem uma corrente de entrada \math{F_{i}} e de saida \math{F_{o}}, com suas devida composições. Nota-se o conjunto de variáveis (coordenadas) \math{\left\lparen T, pH, p\chemical{O_{2}}\right\rparen}, temperatura, pH, e pressão parcial de oxigênio respectivamente, sugerindo uma monitoração dessas variáveis, porque elas são os parametros de interesse para o processo bioquímico apresentado. Portanto, um bioreator pode ser considerado um tanque, ou volume de controle projetado especificamente para adaptar condições otimas para uma cultura de microorganismos metabolisar o alimento presente segundo uma estequimetria ideal, ou proxima disso \cite[authoryear][ref::erickson_0].

  Quando tratamos de processos aeróbicos, é necessário fornecer oxigênio para o desenvolvimento da cultura. Nessa perspectiva, boa parte dos reatores destinados ao cultivo de organismos aeróbicos necessita de agitação para uniformizar a concentração de oxigênio em todo o meio, o que é naturalmente dificultado pela baixa solubilidade do oxigênio em água. Isso pode ser realizado por meio de reatores de mistura, colunas de ar dissolvido, membranas ou pelo aumento da área de interface líquido-gás, o que não ocorre necessariamente em processos anaeróbicos ou em sistemas envolvendo culturas mistas \cite[authoryear][ref::erickson_0]. Essa última opção tem como proposta estimular a competição entre diferentes variantes presentes com base na seleção natural. Portanto, os tipos de cultura influenciam diretamente no projeto físico do biorreator.

  Outra condição de operação que influencia drasticamente tanto o projeto físico quanto o modelo matemático de um reator é a fase dos reagentes e produtos. Em fase líquida, a cinética bioquímica depende da concentração dos reagentes e, em alguns casos, também da concentração dos produtos. Para muitas reações bioquímicas, essa abordagem cinética pode ser tratada, com boa aproximação, considerando dependência de primeira ordem em baixas concentrações e de ordem zero em concentrações elevadas \cite[authoryear][ref::erickson_0]. Em fase sólida, a cinética bioquímica torna-se mais dependente da disposição espacial do substrato sólido e do contato superficial deste substrato com o microrganismo. Dessa forma, uma solubilidade substancial dos reagentes é essencial para garantir adequada difusão no meio e, consequentemente, a formação do produto de interesse \cite[authoryear][ref::bhargav_0]. Em alguns sistemas, utilizam-se bioenzimas para hidrolisar o substrato em reagentes mais solúveis, aumentando a conversão por unidade de tempo.

  O controle das variáveis físicas envolvidas é importante para a predição da resposta do sistema em relação a qualquer perturbação de entrada, bem como para a manipulação do sistema, de forma que os resultados de réplicas experimentais sejam os mais semelhantes possíveis dentro das limitações de sistemas não determinísticos. Entre as variáveis que podem ser medidas em um biorreator estão temperatura, pressão, pH, potência de agitação, fluxo de gases e líquidos, concentração de gases, concentração de biomassa e tensão interfacial \cite[authoryear][ref::erickson_0].
\stopsection


	% \startsection[title={Fermentação}]
	A fermentação ocupa posição de destaque na indústria biotecnológica, assumindo protagonismo, ainda hoje, em rotas produtivas de compostos essenciais para a sociedade \cite[authoryear][ref::metcalfe_0]. Grande parte das vacinas, medicamentos e biocombustíveis é obtida preferencialmente por processos fermentativos, e sua viabilidade econômica decorre justamente da rota empregada.

	Os microrganismos responsáveis pela fermentação podem ser bactérias, leveduras, arqueas, fungos ou células de mamíferos \cite[authoryear][ref::muller_0]. Existe uma ampla variedade de caminhos metabólicos possíveis considerando o conjunto desses organismos. As leveduras, por exemplo, são responsáveis pela produção de etanol, pão, bebidas alcoólicas, entre outros produtos, reforçando a relevância econômica e social desses microrganismos e de suas aplicações industriais.

	A produção de etanol por fermentação é um dos processos biotecnológicos mais importantes da indústria moderna, sendo amplamente utilizada para a obtenção de biocombustíveis e produtos químicos renováveis \cite[authoryear][ref::dashko_0]. Nesse processo, leveduras como \getbuffer[sacc] convertem açúcares provenientes de matérias-primas vegetais, como cana-de-açúcar, milho e biomassa lignocelulósica em etanol e dióxido de carbono sob condições anaeróbicas \cite[authoryear][ref::bhargav_0]. A fermentação alcoólica apresenta grande relevância econômica, especialmente em países como o Brasil, onde o etanol é utilizado em larga escala como combustível automotivo. Além de sua aplicação energética, o etanol também é empregado nas indústrias farmacêutica, alimentícia e química, muitas vezes na forma de gel ou líquida em diversas concentrações para ser utilizado na esterilização.

	\startplacefigure[title={Ilustração de células de microrganismos fermentativos}, reference={fig:cell_illustration}]
		\startcombination[nx=2]
			{\externalfigure[image/yeast_cell.png][height=200pt]}{(A) Levedura}
			{\externalfigure[image/bacterial_cell.png][height=200pt]}{(B) Bactéria}
		\stopcombination
	\stopplacefigure

	Na \in{Figura}[fig:cell_illustration] tem-se (A) uma célula de levedura e (B) a uma célula de bactéria. Nesta, vê-se algumas particularidades que a diferencia da célula de levedura, o que mais se destaca é a disposição do material genético (\chemical{DNA}), espalhado pelo citoplasma, caracterizando esses microrganismos como procariontes e os colocando em contraste com os microrganismos eucariontes de \chemical{DNA} concentrado no núcleo como a levedura. Existem alguns estudos como o de \cite[authoryears][ref::meadows_0] que demonstra o impacto da modificação genética nuclear de células eucariontes na sua performance. A mitocôndria também possui material genético, embora a extensão dos estudos relacionados à sua manipulação seja menor quando comparada à do núcleo. Um ponto importante na compreensão das reações químicas de oxirredução \cite[authoryear][ref::muller_0] é que elas ocorrem no citoplasma ou são coordenadas na mitocôndria quando associadas a processos aeróbicos. A diferenciação do interior de uma célula de microrganismos fermentativos e do exterior destes é importante para o tratamento de modelagens matemáticas dos seus metabolismos. Nestes modelos, as reações que acontecem no citoplasma ou na mitocôndria são balanceadas para que o fluxo molar total destes meios sejam nulos, possibilitando um tratamento homogêneo de sistema linear, e uma solução para os fluxos de todas reações intracelulares, de tal forma que os fluxos molares internos sejam retroalimentados em equações que modelam o balanço molar do exterior da célula, possivelmente por equações diferenciais \cite[authoryear][ref::sainz_0]. 
\stopsection


	% \startsection[title={Fermentação Alcoólica}]
	%ethanol pathways
	Biocombustivel é uma alternativa energética extensamente adotada pela industria de transporte em todo o mundo como substituto ao combustivel fossil (INSERIR DADOS DO BASIL), dentre eles, o etanol desempenha um papel expessivo na demanda energética (DADOS DE CONTRIBUIAÇAO DO ETANOL). A maior parte de todo o biocombustivel produzido e distribuido comercialmente partem de biomassa lignocelulosica, ou seja, existe uma certa dificuldade técnica de pré tratamento desse material para que o biocombustivel seja sintetizado.

	Existem duas rotas principais para produção de biocombustivel a partir de biomassa lignocelulosica, estas são bioquímicas e termoquímicas. Em rotas bioquimicas o pré tratamento é feito em meio ácido ou acalino, ou tambem com explosão a vapor visando a quebra de ligações celolose-hemicelulose-lignina \cite[authoryear][ref::munasinghe_0]. Essa quebra permite a ação enzimatica de hidrólise, culminando na formação de açucares fermentáveis. Em contrapartida, o tratamento da biomassa nas rotas termoquímicas envolve a gaseificação da biomassa em gás de sintese (\chemical{CO} e \chemical{H_{2}}), para posteriormente ser convertido em biocombustivel, usando catalisadores químicos no processo de {\emph Fischer-Tropsch}, ou microorganismos catalisadores no processo conhecido como fermentação de gás de sintese ou {\emph syngas fermentation} \cite[authoryear][ref::munasinghe_0].

	No que se refere à produção de álcool a partir de fontes imediatas de carbono, como glicose ou sacarose, as leveduras desempenham papel fundamental. Em especial, \getbuffer[sacc] representa um grupo extremamente relevante nesse processo. Do ponto de vista evolutivo, essas leveduras desenvolveram essa capacidade inicialmente em resposta ao surgimento de frutas, fontes naturais de açúcares simples, possibilitando a conversão desses açúcares em álcool tanto em condições aeróbicas quanto anaeróbicas \cite[authoryear][ref::dashko_0].

	Em geral, o metabolismo de um microrganismo pode ser modificado por meio da inserção de genes reacionais em estequiometria adequada para maximizar determinado produto ou inibir mecanismos metabólicos específicos. \cite[authoryears][ref::meadows_0] trabalharam na modificação do genoma de variantes de \getbuffer[sacc] e demonstrou a possibilidade de aumentar significativamente a seletividade de caminhos metabólicos de interesse, especificamente na produção de precursores de ácidos isoprenoides utilizados no tratamento da malária. Em situações como essa, estudadas por \cite[author][ref::meadows_0], o metabolismo torna-se diferente daquele naturalmente encontrado em variantes de \getbuffer[sacc]. Entretanto, ao restringir a análise ao metabolismo nativo genérico, caracterizado principalmente pelo consumo de glicose e pela formação de etanol, é possível visualizar o mecanismo representado na \in{Figura}[fig:alcoholic_fermentation].

	\startplacefigure[title={Mecanismo genérico parcial de fermentação alcoólica}, reference={fig:alcoholic_fermentation}]
		\externalfigure[image/alcoholic_fermentation.png][height=200pt]
	\stopplacefigure

	A \in{Figura}[fig:alcoholic_fermentation] mostra a esquematização da rota de produção de etanol na fermentação alcoólica, onde os açúcares, especialmente glicose e sacarose, são metabolizados pelas leveduras em condições predominantemente anaeróbicas. Nesse processo, a glicose é convertida em piruvato por glicólise e, posteriormente, em etanol e dióxido de carbono. Sendo que durante a glicólise há formação de \chemical{ATP} e redução de \chemical{NAD^{+}}, em que este é responsavel pelo transporte de elétrons no ciclo produtivo, enquanto que na formação de etanol a partir de \chemical{Acetaldeido} o \chemical{NADH} é oxidado, ou seja, um ciclo de oxidação-redução do \chemical{NADH}.

	\startplaceformula[reference={eq:alcoholic_fermentation}]
		\startformula
			\startchemicalformula
				\chemical{C_{6}H_{12}O_{6}}
				\chemical{GIVES}
				\chemical{{2}CO_{2}}
				\chemical{PLUS}
				\chemical{{2}C_{2}H_{5}OH}
			\stopchemicalformula
		\stopformula
	\stopplaceformula

	A \in{Equação}[eq:alcoholic_fermentation] coloca a \in{Figura}[fig:alcoholic_fermentation] em uma estequimetria global do processo, sem levar em consideração os ciclos de renegeneração de \chemical{NADH} e \chemical{ATP}. Alem de que ambas representações deixam de fora considerações sbre mecanismos alternativos recorrentes como a produção de {\bf burirato}, ácido acético e butanol e intermediarios equivalentes a \chemical{Acetyl-CoA}. Na \in{Equação}[eq:alcoholic_fermentation], vê-se formação de dióxido de carbono e etanol deterimento do consumo de glicose em uma rota preferêncial de processos produtivos que visam a maximização de conversão em etanol.
\stopsection


	% \startsection[title={Fermentação de hidrogênio}]
	Os combustíveis fósseis possibilitaram avanços significativos para o desenvolvimento da sociedade moderna ao longo dos séculos, contribuindo para a consolidação da industrialização, da expansão dos meios de transporte e da intensificação da globalização. Embora a implementação inicial de métodos de geração energética em larga escala tenha sido fundamental para o desenvolvimento socioeconômico da humanidade, grande parte desses modelos apresentam impactos expressivos sobre a manutenção da estabilidade ambiental e da saúde do planeta a longo prazo. Atualmente, considerando o avanço tecnológico e a viabilidade de implementação de alternativas ao modelo energético predominantemente utilizado \cite[authoryear][ref::nanda_0], torna-se relevante destacar os avanços científicos relacionados a métodos alternativos de geração de energia renovaveis. Dentre elas, as rotas de conversão de resíduos orgânicos por meio de tratamentos termoquímicos e bioquímicos destacam-se como algumas das opções mais promissoras em substituição aos combustíveis fósseis \cite[authoryear][ref::nanda_0].

	O hidrogênio possui elevado potencial para utilização em larga escala como combustível. Nesse contexto, pode ser empregado diretamente como um biocombustível gasoso ou convertido em hidrocarbonetos e eletricidade por meio de células a combustível \cite[authoryear][ref::das_0]. Para fins de combustão, o hidrogênio apresenta valor calorífico inferior de aproximadamente \math{120 \unit{MJ} \unit{kg}^{-1}} \cite[authoryear][ref::sarangi_0], equivalente ao fornecido por cerca de \math{2.75 \unit{kg}} de gasolina \cite[authoryear][ref::nanda_0]. Tais características, associadas ao fato de sua combustão gerar água como único produto, tornam o hidrogênio uma alternativa energeticamente vantajosa, promissora e sustentável.

	As rotas para produção de hidrogênio são diversas, podendo ter como matéria-prima hidrocarbonetos, gás natural, água ou biomassa. Entre os principais processos, destacam-se a reforma de gás natural, decomposição térmica do gás natural, gaseificação de carvão, oxidação parcial de hidrocarbonetos, pirólise, eletrólise, fotólise, termólise e as rotas biológicas \cite[authoryear][ref::das_0]. Essas rotas, predominantemente termoquímicas e eletroquímicas, em sua maioria, requerem elevadas quantidades de energia durante sua operação, e são consideradas ambientalmente pouco sustentáveis. Em contraste, a produção biológica de hidrogênio ocorre sob condições próximas à temperatura e pressão ambientes, apresentando, portanto, menor demanda energética ao longo do processo \cite[authoryear][ref::das_0].

	Dentro das possibilidades de obtenção biológica de hidrogênio, encontram-se três rotas alternativas na natureza, cada uma com características que podem viabilizar ou inviabilizar sua aplicação, a depender das condições de operação e dos substratos utilizados: biofotólise da água, fotodecomposição e fermentação de compostos orgânicos \cite[authoryear][ref::das_0]. Essas são as principais alternativas para a biossíntese de hidrogênio, realizadas em ambientes anaeróbicos facultativos ou obrigatórios. Porém, tanto na biofotólise quanto na fotodecomposição, a formação de \chemical{H_{2}} é parcial ou completamente inibida na presença de oxigênio.

	A biofotólise adapta o processo de fotossíntese de plantas e algas para a geração de hidrogênio, em vez da produção de biomassa baseada em carbono \cite[authoryear][ref::das_0]. O processo ocorre por meio da absorção de luz, que promove a quebra da molécula de água em gás hidrogênio e gás oxigênio \cite[authoryear][ref::das_0]. Embora existam nuances relacionadas ao mecanismo completo da reação, de modo geral a formação de hidrogênio ocorre pela ação das enzimas hidrogenase ou nitrogenase, como ilustrado na \in{Figura}[fig:biophotolysis].

	\startplacefigure[title={Ilustração do mecanismo de biofotólise da água}, reference={fig:biophotolysis}]
		\externalfigure[image/water_photolysis.png][height=200pt]
	\stopplacefigure

	Na \in{Figura}[fig:biophotolysis], observa-se que a enzima (\chemical{Hidrogenase} ou \chemical{Nitrogenase}) desempenha um papel central na redução de prótons provenientes da fotólise da água. Os elétrons liberados nesse processo são transferidos até a enzima por meio da proteína \chemical{Ferredoxina}, na sua forma reduzida \chemical{Fd_{red}}. A oxidação da \chemical{Ferredoxina} pela \chemical{Hidrogenase} constitui a etapa final da cadeia de transferência eletrônica, na qual os \chemical{H^{+}} são reduzidos, resultando na formação de gás hidrogênio (\chemical{H_{2}}). Em seguida, a \chemical{Ferredoxina} é regenerada na forma oxidada (\chemical{Fd_{ox}}), sendo posteriormente novamente reduzida a \chemical{Fd_{red}}, o que permite a continuidade do ciclo redox. Observa-se também a formação de oxigênio molecular na etapa inicial da decomposição da água, conforme ilustrado na \in{Figura}[fig:biophotolysis]. Nesse contexto, o processo apresenta dependência de sistemas enzimáticos para viabilizar a redução de prótons a hidrogênio, caracterizando-o como um sistema indireto de produção de \chemical{H_{2}}. Em contrapartida, pode ocorrer a evolução direta de hidrogênio a partir da fotólise da água, configurando um sistema denominado direto \cite[authoryear][ref::sarangi_0]. Ambos os mecanismos podem ocorrer nesse tipo de rota metabólica; enquanto, o sistema direto apresenta maior simplicidade, o sistema indireto depende da atuação de enzimas específicas.

	A fotodecomposição de compostos orgânicos, também denominada fotofermentação, ocorre na presença de bactérias fotossintéticas \cite[authoryear][ref::sarangi_0], como as do genero \getbuffer[bac]. Esse processo caracteriza-se pela degradação de espécies químicas por ação da luz, com produção de hidrogênio molecular, de maneira análoga à biofotólise. A principal diferença entre os processos reside na natureza da fonte de alimentação e, consequentemente, na rota metabólica global envolvida, embora o princípio de conversão fotobiológica seja comparável. No caso da fotodecomposição, a matéria orgânica atua como substrato energético e doador de elétrons, permitindo maior flexibilidade espectral de absorção luminosa \cite[authoryear][ref::das_0]. Associado ao consumo de matéria orgânica, esse processo apresenta a vantagem de possibilitar a utilização de resíduos orgânicos como fonte de entrada, principalmente carboidratos \cite[authoryear][ref::ntaikou_0].

	\startplacefigure[title={Ilustração do mecanismo de fotodecomposiçao}, reference={fig:photo_fermentation}]
		\externalfigure[image/photo_fermentation.png][height=200pt]
	\stopplacefigure

	A \in{Figura}[fig:photo_fermentation] apresenta semelhanças com a \in{Figura}[fig:biophotolysis], destacando-se, em particular, a presença da \chemical{Ferredoxina}, a qual desempenha função equivalente em ambos os mecanismos. Esses dois sistemas diferem principalmente quanto à fonte de substrato, sendo que, no segundo caso, ácidos orgânicos são preferencialmente utilizados em substituição à água como doadores de elétrons. Outra analogia relevante pode ser estabelecida com o sistema da \in{Figura}[fig:alcoholic_fermentation], no qual o \chemical{NADH} desempenha papel funcional análogo ao da \chemical{Ferredoxina} nos processos de fotodecomposição e biofotólise da água. Em ambos os casos, esses cofatores atuam como transportadores de elétrons, viabilizando reações de redução na estequiometria final dos processos \cite[authoryear][ref::sarangi_0]. A \in{Equaçao}[eq:photo_fermentation] apresenta uma representação estequiométrica genérica do processo de fotodecomposição.

	\startplaceformula[reference={eq:photo_fermentation}]
		\startformula
			\startchemicalformula
				\chemical{{2}CH_{3}COOH} 
				\chemical{PLUS} 
				\chemical{H_{2}O} 
				\chemical{GIVES} 
				\chemical{CO_{2}} 
				\chemical{PLUS} 
				\chemical{{4}H_{2}}
			\stopchemicalformula
		\stopformula
	\stopplaceformula

	Na \in{Equação}[eq:photo_fermentation], observa-se a conversão do \chemical{Ácido acético} em \chemical{H_{2}}. Contudo, essa representação não explicita a dependência do processo em relação à enzima \chemical{Nitrogenase}, bem como à necessidade de \chemical{ATP} para sua atividade catalítica \cite[authoryear][ref::muller_0].

	A lista extendida das variações de substrato aceitas por sistemas de fotodecomposição de compostos orgânicos inclui monóxido de carbono em conjunto com água, em um processo denominado translação água-gás \cite[authoryear][ref::das_0], onde a estequiometria correspondente é apresentada na \in{Equação}[eq:shift].

	\startplaceformula[reference={eq:shift}]
		\startformula
			\startchemicalformula
				\chemical{CO} 
				\chemical{PLUS} 
				\chemical{H_{2}O} 
				\chemical{GIVES} 
				\chemical{CO_{2}} 
				\chemical{PLUS} 
				\chemical{H_{2}}
			\stopchemicalformula
		\stopformula
	\stopplaceformula

	A \in{Equação}[eq:shift] solidifica o conceito nomeado translação água-gás, ou {\emph water-gas shift}, que se encaixa no contexto da movimentação do oxigênio da água para o monóxido de carbono, oxidando o carbono, reduzindo o hidrogênio, e formando gás hidrogênio.

	Para ambos os sistemas envolvendo fotodecomposição, não há a formação de oxigênio molecular no meio reacional. Dessa forma, o planejamento do sistema é simplificado, uma vez que não há necessidade de controle rigoroso da presença de \chemical{O_2}, o qual pode atuar como um forte inibidor da atividade de hidrogenases. O efeito inibitório do oxigênio está associado à sensibilidade dessas enzimas ao meio oxidante, podendo ocorrer inativação reversível ou irreversível do sítio ativo contendo metais de transição, particularmente centros \chemical{Fe-Fe} ou \chemical{Ni-Fe}. Nessas condições, a presença de \chemical{O_{2}} leva à oxidação desses centros catalíticos, reduzindo significativamente a eficiência da redução de prótons a \chemical{H_{2}} (CITAÇAO). Em sistemas como a biofotólise da água, essa limitação exige estratégias adicionais de separação temporal ou espacial entre a etapa de geração de oxigênio e a etapa de produção de hidrogênio, a fim de preservar a atividade enzimática.

	A fermentação de hidrogênio constitui o último dos três principais métodos biológicos de produção de \chemical{H_{2}}, apresentando cinética e condições reacionais distintas, o que resulta em características sistemáticas particulares. Bactérias fermentativas podem apresentar elevadas taxas de evolução de hidrogênio, com capacidade de produção tanto na presença quanto na ausência de luz, não havendo dependência de iluminação, o que permite sua atividade contínua em diferentes ciclos ambientais. Além da produção de \chemical{H_{2}} em quantidades relativamente elevadas quando comparadas a outros sistemas biológicos, esses microrganismos geralmente apresentam taxas de crescimento adequadas para aplicações em processos produtivos, o que pode favorecer o aumento da conversão sob condições ambientais otimas \cite[authoryear][ref::das_0]. Também denominada fermentação escura, esse processo pode ocorrer em condições estritamente anaeróbias ou facultativamente anaeróbias, dependendo do microrganismo envolvido, com uma estequiometria genérica apresentada na \in{Equação}[eq:dark_fermentation]. Exemplos incluem bactérias do gênero \getbuffer[clos] \cite[authoryear][ref::muller_0].

	\startplaceformula[reference={eq:dark_fermentation}]
		\startformula
			\startchemicalformula
				\chemical{C_{6}H_{12}O_{6}}
				\chemical{PLUS}
				\chemical{{2}H_{2}O}
				\chemical{GIVES}
				\chemical{{2}CH_{3}COOH}
				\chemical{PLUS}
				\chemical{CO_{2}}
				\chemical{PLUS}
				\chemical{{4}H_{2}}
			\stopchemicalformula
		\stopformula
	\stopplaceformula

	A \in{Equação}[eq:dark_fermentation] exemplifica uma das possíveis reações globais realizadas por microrganismos fermentativos anaeróbios heterotróficos, na qual ocorre o consumo de \chemical{Glicose} com formação de \chemical{ácido acético} e \chemical{H_{2}}. O \chemical{ácido acético} é um ácido orgânico que pode ser posteriormente consumido por microrganismos fotodecompositores, de modo que a produção global de hidrogênio é potencializada pela combinação entre fermentação escura e fotofermentação \cite[authoryear][ref::muller_0]. Dessa forma, um arranjo em série das reações representadas nas \in{Equação}[eq:dark_fermentation] e \in{Equação}[eq:photo_fermentation] permite maior aproveitamento dos ácidos orgânicos intermediários, resultando em conversão globais mais eficientes de substrato em \chemical{H_{2}}. Outro arranjo possivel é o reaproveitamento dos subprodutos da fermentação escura na eletrólise microbiana. Esta possui uma estequimetria semelhante a fotofermentação, com distinções em torno da força motriz, mudando para oxiredução em eletrodos.
\stopsection


	% \startsection[title={Sistema de controle}]
  Controlar o curso do processo fermentativo em um bioreator é essencial para o sucesso de um procedimento passivel de ser replicado ou otimizado. Isso pode ser realizado somente com o auxilio de um sistema de controle robusto, onde variáveis como temperatura, pH ou concentracão de alguma espécie química precisam ser monitoradas e modificadas direta ou indiretamente sob o pretexto de manutenção de um ambiente propício ao crescimento da cultura ou consumo de substrato.

  O ponto de operação ou {\emph set point} de cada uma das variáveis controladas, tem de ser definido para o experimento, ou processo conduzido em dado momento. Mas para atingir o objetivo primeiro de monitorar essas variáveis, tem-se de adiquirir sensores específicos e, em seguida, integra-los no sistema de controle.

  \startplacefigure[title={Esquematico de um bioreator fermentativo}, reference={fig:bioreactor_control}]
    \externalfigure[image/bioreactor_control.png][height=200pt]
  \stopplacefigure

  Um digrama de processo é mostrado na \in{Figura}[fig:bioreactor_control], em que ve-se a direita o controlador recebendo sinais dos sensores de temperatura \math{T}, espuma, \chemical{pH}, \chemical{O_{2}} e vazão de entrada de ar \math{Q_{air}}, e respondendo segundo a lógica de funcionamento do sistema em torno dos {\emph set points} para as valvulas, bombas peristalticas, aquecedor e rotor. A entrada de ar é responsável pela alimentação de gases, importantes para o metabolismo dos microorganismos, \chemical{O_{2}} para fermentação aeróbica. O controle do \chemical{pH} acontece sob atuação da bomba peritaltica interconectadas com resenvatórios de ácido e base, em função do desvio de ácidez do sistema interno, assim como o antiespuma, que é injetado em resposta ao aumento de espuma. Alimentação de substrato pela bomba peristaltica, acontece segundo a politica de alimentação do experimento, podendo ser intercalda, ou única.

  Ainda na \in{Figura}[fig:bioreactor_control] a temperatura interna pode ser controlada a partir do fluxo de água, que entra resfriada pronta para reduzir a temperatura interna do bioreator, ou disponível para aquecimento quando circula pelo aquecedor ativado segundo o comando do controlador, a depender da temperatura interna.

  Cada variante de microorganismo fermentador possui uma região de temperatura preferida para a manutenção das suas atividades metabolicas, ou mesmo, do seu crescimento e existencia no broto, portanto, tratam-se variantes em um espectro de temperatura. Elas podem ser classificadas como mesofílicas, termofílicas ou termofílicas extremas, sendo cada uma dessas denominações responsável pela representatividade da faixa de temperatura obrigatória da variante.

  {\emph \bf Gráfico do comportamento inibitorio da temperatura na conversão para um dada cultura pura}

  Além da alta sensibilidade a temperatura, os microorganismos fermentativos tambem possuem sensibilidade ao pH. A depender do intervalo de pH do experimento, pode não haver nenhuma atividade metabolica voltada para produção de algum composto, por parte de algumas variantes, ou seja, caracterizando inibição de rotas, ou de todo metabolismo dessas variantes.

  {\emph \bf Gráfico do comportamento inibitorio do pH na conversão para um dada cultura pura}
\stopsection


	\startsection[title={Modelagem matemática}]
  \writestatus{INFO}{Text width = \the\hsize}
  Monitorar e controlar processo industriais é em certa medida a principal atividade dos sensores físicos. Porem, devido a grande quantidade de informação que esses sensores conseguem carregar para fora do sistema, é possivel levar esses dados em consideração em interpretações mais abstratas do processo. A construção de um modelo preditivo de qualidae do processo é uma das possibilidades contempladas pela interpretação do contexto geral dado pelos sensores físicos, esses modelos são denominados \getbuffer[sens]. Este termo é derivado da combinação entre as palavras {\emph software} e {\emph sensor}, dado o sentido da existencia de um sensor pautado em uma lógica programada ao inves da sensibilidade de um material partindo diretamente de um princípio físico.

  {\emph Soft sensors} podem ser distinguidos em duas classes, a primeira é centrada em dados, nomeada {\emph data-driven} e a segunda formada por uma descrição completa de um modelo físico-químico, denominada {\emph model-driven}. Estas duas classes são fundamentalmente distintas, podendo ser situadas nos extremos de um espectro de implementação fenomenologico do sistema. {\emph Data-driven} foca sua capacidade computacional em entender a correlação dos dados de entrada com o termo de qualidade na saida so sistema, enquanto aplicações {\emph model-driven} necessita de uma modelagem dos princípios físicos e químicos acontecendo no plano de fundo do processo, e sua correlação matemática com os parâmetros de qualidade de saida do sistema. Consequentemente, a aplicação de {\emph model-driven} \getbuffer[sens] é tecnicamente dificultada por definição, fazendo com que a maior parte das implementações sejam voltadas para modelos pautados em variáveis de processos de entrada. Portanto, neste trabalho, quando referenciado \getbuffer[sens] lê-se {\emph data-driven} \getbuffer[sens], exceto quando devidamente especificado.

  {\emph \bf Espectro dos modelos de soft sensors entre model-driven e data-driven comparando com modelos white-box, gray-box e black-box}

  Modelagem de \getbuffer[sens] tem por objetivo agregar variáveis de entrada, juntamente com o potêncial dessas variáveis de carregar informações sobre o processo, e devolver na saída características de qualidade do processo \cite[authoryear][ref::yang_0]. Ou tambem sinalizar desvios de idealidade, informando ao operador problemas ocorridos no processo, possibilitando uma detecção automatizada de problemas difíceis de serem encontrados \cite[authoryear][ref::kadlec_0]. Esta ultima prerrogativa é especialmente util quando trata-se de processos continuos, nos quais o reator encontra-se idealmente em regime permanente, porem dada a presença de microorganismos há imprevisibilidade quanto a resposta aos minímos estimulos que fogem da idealidade, ou mesmo uma resposta átipica a condições ideais de operação devido ao estress acumulado pela cultura. Com um modelo treinado em especificidades como essa, é possivel que haja uma predição da problematica antes memos que o curso divergente seja capaz de causar maiores danos na produtividade ou na composição de saida.

  Sobre os problemas encontrados em implementações de \getbuffer[sens] pode-se destacar dificuldades com tratamento de medições com ruido, ou fora do padrão e colinearidade entre variáveis \cite[authoyear][ref::kadlec_0]. Suas resoluções são geralmente tratadas com interferência manual sobre os parametros, modelo e coordenadas utilizadas no desenvolvimento da aplicação, a fim de remover pontos fora da curva esperada e como tambem excluir interdependência entre variáveis de entrada. Essas soluções se dão principalmente devido ao entendimento do processo, e dificilmente poderiam ser de outra forma (CITAÇAO).

  O formato dos dados esperados para a contrução dos modelos de \getbuffer[sens] são separados em conjunto de entrada \math{X} e conjunto alvo \math{Y}, formando o conjunto de dados \math{S}.

  \startplaceformula[reference={eq:data_set}]
    \startformula
      S = \left\lbrace \left\lparen X, Y \right\rparen \right\rbrace
    \stopformula
  \stopplaceformula

  Na \in{Equação}[eq:data_set] observa-se os conjuntos mencionados, dos quais as iterações programaticas do algoritimo particular tem por objetivo aproximar uma função alvo \math{f} que transforma \math{X} em \math{Y}, da forma que aparece na \in{Equação}[eq:data_set_train]. A dimensão vetorial dos elementos de \math{X} e \math{Y}, vão depender da quantidade de variáveis medidas em cada uma das amostragems. Ou seja, \math{\vec{x}} pode ser um elemento de \math{X} com dimensão \math{m} e \math{\vec{y}} um elemento de \math{Y} com dimensão \math{\gamma}. Os elementos do conjunto de entrada \math{X} e do conjunto de saida \math{Y} não possuem qualquer correlação dimensional, porem esses conjuntos compartilham a necessidade de possuir a mesma quantidade de elementos para que o espaço amostral seja valido, visando a formação de pares \math{(\vec{x}, \vec{y}) \in S} para definição da função \math{f}.

  \startplaceformula[reference={eq:data_set_train}]
    \startformula
      f : X \rightarrow Y
    \stopformula
  \stopplaceformula

  Lê-se na \in{Equação}[eq:data_set_train] uma aplicação \math{f} que leva do conjunto \math{X} ao conjunto \math{Y}. Percebe-se implicitamente pela definição apresentada que \math{f} é desconhecida inicialmente, e toda a teoria desenvolvida em torno de \getbuffer[sens] tem por objetivo encontra-la, mesmo que por trás de uma caixa preta.

  \startitemize
    \starthead{Multiple linear regression} 
      Este modelo é considerado estatístico classico. Baseia-se em um método de ajuste de coeficientes para um função linear pré-definida da forma da \in{Equação}[eq:regresssion], pautada nos principios de minimização de erros quadrados.

      \startplaceformula[reference={eq:regresssion}]
        \startformula
          y = \beta_{0} + \sum_{j=1}^{n} \beta_{j} x_{i j}
        \stopformula
      \stopplaceformula

    \stophead

    \starthead{multilayer perceptron}
    \stophead

    \starthead{partial least square}
    \stophead

    \starthead{recurrent neural networks}
    \stophead

    \starthead{long-term and short-term memory}
    \stophead
  \stopitemize
\stopsection


	\startsection[title={Modelo fermentativo}]
	A modelagem do processo de fermentação é essencial para seu estudo, porque toda otimização do processo é tão bem feita quanto o modelo físico-químico permite \cite[authoryear][ref::ureta_0]. Seguindo as diretrizes do modelo paramétrico formulado, é possivel estimar os valores desses parâmetros a partir do processamento de dados de um conjunto de experimentos cuidadosamentes planejados e executados. Validação desses resultados é uma ação rotineira no exercício de desenvolvimento de modelos devidamente equacionados \cite[authoryear][ref::ureta_0].

	O primeiro passo para o desenvolvimento de um modelo fermentativo é a analise qualitativa da estrutura do sistema fermentativo, e entendimento das rotas metabólicas. Em seguida, deve-se formular um modelo matemático que segue as definições a priori utilizadas na etapa anterior. Por último é necessário a identificação dos valores numéricos que substituirão os parâmetros e constantes do modelo, baseados em dados experimentais reais de um processo.

	Criar um modelo matemático envolve a etapa de determinação da estequiometria de reação e a etapa de formulação da taxa de reação. A etapa estequiométrica tem suas particularidades e pode englobar diversos mecanismos consideravelmente complexos, consequentemente, a taxa de reação \math{r}, que tem uma relação causal com a estequiometria, geralmente acompanha sua complexidade, acarretando em um uma cinética extensa. Mas existem modelos que tratam de generalizar comportamentos de determinadas reações, como mostra o \in{Sistema de Equações}[eq:rates], onde a mais simples tem a forma da reação elementar de ordem \math{1} no substrato \math{S}.

	\startplaceformula[reference={eq:rates}]
		\startformula
			r_{1} = k S \breakhere
			r_{2} = k S^{n} \breakhere
			r_{3} = \frac{k S}{K + S} \breakhere
			r_{4} = \frac{k S}{K + S} I
		\stopformula
	\stopplaceformula

	Em que \math{K} é a constante limitante da taxa de reação, \math{k} a constante de taxa de reação, \math{I} representa a inibição, e \math{S} a concentração de algum substrato ou biomassa.

	{\emph Anaerobic Digestion Model no 1} (ADM1) é uma das alternativas para modelos fermentativos praticos de aplicação industrial, que envolve diversos componentes, tipos de microorganismos, reações e inibições com base em equações diferenciais \cite[authoryear][ref::batstone_0]. Neste modelo é possivel representar as diferentes etapas da vida de um microorganismo atraves da alteração das constantes de reação, permitindo modelar, fases de crestimento, morte e atraso, para isso é necessário introduzir equações diferenciais que respeitam as variações de biomassa como da \in{Equação}[eq:biomass_differential].

	\startplaceformula[reference={eq:biomass_differential}]
		\startformula
			\frac{\dd X_{a}}{\dd t} = k_{a1} X_{a} - k_{a2} X_{a}
		\stopformula
	\stopplaceformula

	Em que \math{X} é a concentração de biomassa ativa, e \math{k_{a1}, k_{a2}} são as constantes de crescimento e desativação respectivamente.

	Os mecanismos de reação modelados no ADM1 começam com a desintegração de um particulado complexo em carboidrato, proteina, gordura, inerte soluvel e inerte insoluvel. A partir deste ponto, acontece a hidrolise desses compostos em monosacarídeos, amino ácidos e ácidos graxos, seguindo para a formação de compostos como \chemical{valerato}, \chemical{butirato}, \chemical{propianato}, \chemical{acetato}, \chemical{hidrogênio} \chemical{dióxido de carbono}, \chemical{metano}, entre outros. Estas etapas, mostradas na \in{Figura}[fig:desintegration_path], o tornam um modelo genérico capaz de representar uma cultura mista de biomassa acídogênica, acetogênica e metanogênica, com boa aproximação para sistemas industriais.

	\startplacefigure[title={Caminho de desintegração e hidrolise do ADM1}, reference={fig:desintegration_path}]
		\externalfigure[image/desintegration_path.png][height=200pt]
	\stopplacefigure

	A forma basica de uma equação diferencial para líquidos no modelo é dado na \in{Equação}[eq:soluble_differential].

	\startplaceformula[reference={eq:soluble_differential}]
		\startformula
			\frac{\dd S}{\dd t} = \frac{Q}{V} \left\lparen S_{i} - S_{o} \right\rparen + \sum_{j} \rho_{j} \nu_{ji}
		\stopformula
	\stopplaceformula

	Em que vazão, volume, concentração de entrada do solúvel, concentração de saida do solúvel, taxa cinética e coeficiente cinético são representados por \math{Q}, \math{V}, \math{S_{i}}, \math{S_{o}}, \math{rho_{j}} e \math{\nu_{ji}} respectivamente.

	Outro equacionamento importante do ADM1 é o equilíbrio ácido-base, necessário para o calculo do pH, da possivel inibição de taxas de reações e do crescimento da biomassa. Este equacionamento é dado pela utilização da concentração dos ions na solução, e das constantes de equilíbrio previamente definidas, culminando em uma EDO, ou em uma equação algebrica diferencial de solução iterativa.
\stopsection

\stopchapter
